\section{Introduction}\label{sec:introduction}

Cardy's formula predicts that a critical two-dimensional percolation crossing
probability is determined by a conformal parameter of the marked domain
\cite{Cardy1992}. Smirnov proved the formula for site percolation on the
triangular lattice \cite{Smirnov2001}. Schramm's Problem~2.12 asks for the
corresponding conformal-invariance theorem for Voronoi percolation
\cite{Schramm2006}. We prove the annealed Cardy-crossing aspect of this
conjecture for the Poisson model, with no regularity assumption on the
boundary beyond the Jordan property.

\subsection{The crossing law}

Identify the oriented plane with $\C$. For $\eps>0$, let $\eta=X_\eps$ be
a homogeneous Poisson point process of intensity $\eps^{-2}$. For $v\in\eta$
its Voronoi cell is the closed set
\[
 \mathcal V_v=\{z\in\C:|z-v|\le |z-u|\text{ for every }u\in\eta\}.
\]
Independently of $\eta$, color its sites black or white with probability
$1/2$ each, and give each cell the color of its site. Let $\mathcal B_\eps$
be the union of the closed black cells. Probabilities and expectations under
the joint law of the points and colors are denoted by $\PP$ and $\EE$.
Conditional probabilities given the points are denoted by $\Peta$;
$\Esp$ denotes averaging over independent fair colors on a fixed graph.
Thus all probabilities in the main theorem are annealed.

Let $D$ be a bounded Jordan domain and let $a,b,c,d$ be distinct points of
$\partial D$ in counterclockwise order. Boundary arcs always include their
endpoints. Define $\mathcal C_\eps(D;a,b,c,d)$ to be the event that
$\mathcal B_\eps\cap\overline D$ contains a connected subset meeting both
the arc $ab$ and the arc $cd$. The conformal map from $D$ to
$\mathbb H=\{z:\operatorname{Im}z>0\}$ which sends $a,c,d$ to
$0,1,\infty$ extends to the boundary in the sense of Carath\'eodory
\cite[Section~2.1, p.~18]{Pommerenke1992}.
Its value at $b$ is a uniquely determined number $x\in(0,1)$.

\begin{theorem}[Annealed Voronoi crossing law]\label{thm:cardy}
For every bounded Jordan domain $D$ and every four distinct
counterclockwise boundary marks $a,b,c,d$, one has
\begin{equation}\label{eq:cardy}
 \lim_{\eps\downarrow0}\PP\bigl(\mathcal C_\eps(D;a,b,c,d)\bigr)
 =F(x)
 :=\frac{\displaystyle\int_0^x[u(1-u)]^{-2/3}\,\dd u}
          {\displaystyle\int_0^1[u(1-u)]^{-2/3}\,\dd u}
 =\frac{\Gamma(2/3)x^{1/3}}
        {\Gamma(1/3)\Gamma(4/3)}
        {}_2F_1\!\left(\frac13,\frac23;\frac43;x\right).
\end{equation}
\end{theorem}

The theorem resolves positively the annealed Cardy-crossing conjecture
for critical Poisson--Voronoi percolation, the crossing-probability aspect
of Schramm's Problem~2.12 above. Its scope is a limit of crossing
probabilities averaged over geometry and colors. The statement does not
assert convergence of the full interface law or a quenched crossing limit
for arbitrary Jordan domains.
In particular, the result concerns the specified closed-cell crossing
event in every fixed Jordan domain, including domains for which a smooth
boundary approximation alone would not immediately identify that event.

For a fixed rectangle, the established concentration of conditional
crossing probabilities also gives a quenched $L^2$ limit; see
Corollary~\ref{cor:quenched-rectangles} below.

\subsection{Prior work}

The conformal-invariance conjecture for critical crossings developed from
numerical and universality investigations. Langlands, Pouliot and
Saint-Aubin explicitly credit Aizenman with the conformal formulation
\cite[Sections~2.1 and~2.4]{LanglandsPouliotSaintAubin1994}.
Cardy's conformal-field-theory calculation identifies the predicted
crossing function \cite{Cardy1992}. Smirnov's proof on the triangular
lattice uses color-switching identities to construct conformal
observables \cite{Smirnov2001,Smirnov2001Expanded}. Carleson's
representation of Cardy's formula as a linear boundary value in an
equilateral triangle, explained in \cite[Corollary~3]{Smirnov2001Expanded},
also describes the boundary problem at the end of our argument.

The same Cardy function is obtained for critical square-lattice bond
percolation under uniform marked-Jordan approximation
\cite[Theorem~1.2]{OpenAISquareBond2026}.
There the event is an actual open-edge path with all boundary edges
sampled independently and no external connections; here it is an annealed
closed-cell crossing in a fixed Jordan domain. The square-bond theorem
is not used in our proof.

Voronoi percolation provides a random planar geometry whose law has no
preferred lattice direction. Benjamini and Schramm proved a distinct
high-intensity invariance under smooth conformal changes of metric,
with the sampling measure held fixed \cite{BenjaminiSchramm1998}.
Bollob\'as and Riordan established that the critical probability of the
planar Poisson model is $1/2$ \cite{BollobasRiordan2006}.
Uniform box-crossing estimates were established by Tassion \cite{Tassion2016};
quenched crossing estimates were developed by
Ahlberg, Griffiths, Morris and Tassion \cite{AGMT2016}.
Vanneuville's quantitative work supplies comparisons between quenched
and annealed arm probabilities, as well as the crossing estimates used
here \cite{Vanneuville2025}. The conformal-invariance conjecture and its
relation to these partial results are discussed in
\cite[Section~1.2]{Vanneuville2021}. These estimates control crossings and
exceptional geometry; identifying the limiting function requires an
additional argument. Proposition~\ref{prop:probability-input} states the
established estimates at the precise hypotheses used below.

The challenge is to obtain an analytic relation that survives the
irregular embedding. In a related setting, Beffara separates a
color-switching identity from the geometric defect of a non-equilateral
embedding \cite[Sections~3.2--3.3, equations~(13)--(14)]{Beffara2014Mesoscopic}.
His crossing theorem concerns periodic triangulations subdivided into
triangular-lattice pieces, under uniform box-crossing and a specified
relation between the two scales; its limit agrees with triangular-lattice
percolation up to a real-linear transformation determined by the initial
embedding. Here the geometry is Poisson--Delaunay;
we recover the needed relation by transporting exact closed-circuit
identities to open portions of random boundaries.

\subsection{From circuit identities to the crossing law}

The proof begins with the triangular graph joining sites of neighboring
Voronoi cells and its dual graph. Dual edges separating opposite colors
form ordinary cluster-boundary circuits. For each such circuit $L$, we
delete all points of $L$, including its vertices, and retain each
connected component $C$ of the remaining dual graph. An unused edge whose two endpoints lie
on $L$ then leaves an isolated open-edge component. The label is the
pair $(L,C)$, not just the
geometric side of $L$. Except for isolated open-edge components, these
components have simple monochromatic primal boundaries $P$, called
\emph{rims}, oriented with their component on the left. Keeping the full
component label, including isolated edges, is what makes the exact
orthogonality identity of Section~\ref{sec:fields} hold.

Each triangle is the image of a reference equilateral triangle under
an orientation-preserving real-affine map. The complex-linear and
antilinear parts of this map split each clockwise side vector into two
terms. For a smooth test function $f$, take the contributions for the
side belonging to $P$ in each triangle immediately to its right, and
weight them by $f$ at the triangle's barycenter. Summing the complex-linear
terms gives $U_P(f)$; summing the antilinear terms gives $Q_P(f)$.
Thus $Q$ records the departure from equilateral geometry, while
$V_P(f)=U_P(f)+Q_P(f)$ is the weighted tangent displacement. Orthogonality gives a separate exact second-moment
identity for each Fourier field when summed over all component labels,
including isolated open edges.
The complete tangent sum is $V_P(1)=0$ on every closed rim.
Stationarity makes the corresponding macroscopic closed-rim sums of $Q$
small. The main task is to open that identity: along a common microscopic
subsequence, we prove
\[
                    Q_P(f)-aV_P(f)\longrightarrow0
\]
in probability, uniformly over rims extending beyond a fixed collar of
$\supp f$. The scalar $a$ is independent of the later test and bulk patch;
Theorem~\ref{thm:bulk-transport} states the precise local formulation.

For this passage to open arcs, we retain local colored configurations near
endpoints in deterministic regions called \emph{plugs}. Separated
two-color corridors join them into closed experiments. Surrounding circuits restrict the path to the same two contacts, and the
absence of an alternative monochromatic path on its right forces agreement
when the exterior configuration is changed. For each fixed connector, conditional resampling produces a function of
the two endpoint states, common to its completions, whose sum around each
fixed cycle is small.
Comparing closures and inserting intermediate plugs decomposes this
function into an affine displacement, two local endpoint potentials, and
a controlled remainder. Sections~\ref{sec:plugs} and~\ref{sec:affine}
construct this endpoint transport before it is used on a random arc.

Two issues remain in returning to actual rims. The local geometry and
colors have already selected the arc, so the comparison must retain their
conditioning. At each fixed comparison scale, arm estimates provide a
finite menu of plug models densely along every relevant rim and a finite
deterministic family of connecting patterns. The comparison is made with
unnormalized subprobabilities dominated by the original law; there is no
division by a configuration-dependent fitting probability. A second issue
is cancellation: the fields are signed, so a small change in their
coefficients is not controlled by a path-length estimate. Deterministic
spatial refinement, a near-return bound, and the orthogonality identity
supply that control. Sections~\ref{sec:pins} and~\ref{sec:actual-arcs}
complete these two steps in the order needed for bulk transport.

The final steps prove $1-2\operatorname{Re}a>0$ for the scalar $a$
appearing in the transport relation. In particular, the coefficient
$1-a$ needed by the observable is nonzero. Complete a finite discrete
disk by three exterior marked vertices; a component observable at this
closing triangle records which boundary connection occurs. One of its
boundary coefficients is exactly the crossing probability. Orthogonality
and the transport relation give the contour cancellation needed for a
holomorphic limit. The scalar $a$ need not be identified: once
$1-a\ne0$, it disappears from the limiting Cauchy--Riemann equation.
The resulting boundary-value problem has the conformal map to an
equilateral triangle as its unique solution
(Propositions~\ref{prop:observable} and~\ref{prop:holomorphy}, and
Lemma~\ref{lemma:boundary-uniqueness}). This common boundary problem
identifies every subsequential crossing limit.
A deterministic comparison at fixed buffers in a Schoenflies chart then
transfers this crossing law to the exact closed-cell event in arbitrary
Jordan domains.

\subsection{A quenched consequence for rectangles}

\begin{corollary}[Quenched Cardy limit for fixed rectangles]\label{cor:quenched-rectangles}
Let $D$ be a fixed bounded nondegenerate open rectangle and let $a,b,c,d$
be its corners in counterclockwise order, with $ab$ and $cd$ the chosen
opposite sides. Write $x_D$ for its conformal parameter above and set
\[
 Q_\eps(\eta)=\Peta\bigl(\mathcal C_\eps(D;a,b,c,d)\bigr),
 \qquad \eta=X_\eps.
\]
If $\PP_{\mathrm{geom},\eps}$ and $\EE_{\mathrm{geom},\eps}$ denote
probability and expectation for the Poisson geometry at mesh $\eps$, then
\[
 \lim_{\eps\downarrow0}
 \EE_{\mathrm{geom},\eps}\bigl[|Q_\eps-F(x_D)|^2\bigr]=0.
\]
In particular, for every $t>0$,
$\PP_{\mathrm{geom},\eps}(|Q_\eps-F(x_D)|>t)\to0$.
These limits use the geometry law at each mesh, without a coupling across
meshes.
\end{corollary}

\begin{proof}
Almost surely only finitely many Voronoi cells meet $\overline D$, and
each cell clipped to $\overline D$ is closed and convex. Every connected
component of a finite union of such clipped cells is path connected, by
concatenating segments through a chain of intersecting cells. Thus, almost
surely in the geometry and for every coloring, our connected-set crossing
event agrees with the continuous closed-cell crossing event in
\cite[Definition~1.1]{Vanneuville2025}.

After an orientation-preserving isometry, write
$\overline D=[-w,w]\times[-h,h]$, with the chosen sides vertical.
Scaling by $\eps^{-1}$ sends the full-plane intensity-$\eps^{-2}$ Poisson
process to unit intensity and this rectangle to
$[-\lambda R,\lambda R]\times[-R,R]$, where $\lambda=w/h$ and
$R=h/\eps$. The Voronoi cells and their conditional fair coloring
commute with this similarity. Therefore
\cite[Theorem~1.3]{Vanneuville2025} gives an absolute $\alpha>0$ and a
constant $C_\lambda<\infty$ such that, for $0<\eps\le h$,
\[
 \Var_{\mathrm{geom},\eps}(Q_\eps)
 \le C_\lambda(\eps/h)^\alpha.
\]
Let $m_\eps=\PP(\mathcal C_\eps(D;a,b,c,d))$. Conditioning on the
geometry gives $\EE_{\mathrm{geom},\eps}Q_\eps=m_\eps$, while
Theorem~\ref{thm:cardy} gives $m_\eps\to F(x_D)$. Hence
\[
 \EE_{\mathrm{geom},\eps}|Q_\eps-F(x_D)|^2
 =\Var_{\mathrm{geom},\eps}(Q_\eps)+|m_\eps-F(x_D)|^2
 \le C_\lambda(\eps/h)^\alpha+|m_\eps-F(x_D)|^2
 \longrightarrow0.
\]
Markov's inequality gives the assertion in probability.
\end{proof}

\subsection{Geometry, bulk patches and notation}

Write $\omega=\exp(2\pi\ii/3)$. Unless specified otherwise, sites, edges
and triangles belong to the primal Delaunay triangulation. Either
color can be the path color; we describe it as black when making a
construction. Distances, diameters and balls are Euclidean, and complex
vectors use the positive orientation of the plane.

We recall the geometric facts behind this convention. Almost surely,
the Poisson process is locally finite and has neither collinear triples
nor cocircular quadruples, by diffuseness and a countable exhaustion.
For a finite set in general position, empty-circumdisk triangles form
the triangulation of its convex hull: lift the sites to the paraboloid
and project the lower faces. Their circumcenters are Voronoi vertices,
and shared cell sides correspond to primal edges. The local
determination estimate in Lemma~\ref{lemma:mesh} makes these finite
constructions stabilize on every compact set along an exhaustion.
It also bounds the cells of sites in any fixed compact set. Thus the
full-plane cells form a locally finite polygonal tiling with trivalent
meeting points. Sites of touching cells are primal neighbors, including
when the cells meet at a vertex. The primal graph can equivalently be
drawn through the tiling by joining sites to interior points of shared
cell sides. This drawing has the same planar orders as the straight
Delaunay drawing.

\begin{definition}[Bulk setting]\label{def:bulk-setting}
A bulk setting is either the full-plane sample or a finite simple
triangulation of the sphere carrying independent fair site colors
conditional on its graph and geometry. In the latter case an oriented
open planar patch is coupled to a full-plane sample so that, on every
fixed compact subset of the patch, the graphs and their colors agree
with probability tending to one as $\eps\downarrow0$. The graph and
chart agreement event is independent of color sampling conditional on
the geometry. Spatial weights are used only on triangles in the
agreement patch. Compactly supported tests may be set to zero on the
exceptional geometry event.
\end{definition}

The finite completions will be useful when the observable is formed
near the boundary. No geometric control of their triangles outside
the patch is required for the local field identities.

We use the Jordan--Schoenflies theorem in its prescribed-boundary form
\cite[Theorem~10.4]{Moise1977}, its extension to a simple planar arc
\cite[Premi\`ere Partie, Chapitre~I, Section~17]{Antoine1921},
and Carath\'eodory's boundary extension theorem. We also use
the following form of the Rad\'o convergence theorem
\cite[Section~1, pp.~182--183]{Rado1923}. If Jordan
boundaries have parametrizations converging uniformly and in order to
a Jordan boundary, their Riemann maps, normalized at a common interior
point and with positive derivative, converge uniformly in the
closed-disk parametrization to the limiting chart. When a prescribed
ordered cut is to be realized by an ambient change, we give the
extension argument explicitly in Lemma~\ref{lemma:exact-fit}.
Paths, band sides and gates are always considered together with their
planar embeddings.

\subsection{Reading the proof and the order of limits}

Section~\ref{sec:fields} establishes the exact fields and the stationary
closed-sum estimate. Section~\ref{sec:probability} states the percolation
inputs and derives the local arm and mesh estimates.
Sections~\ref{sec:plugs} and~\ref{sec:affine} construct the state
transport and its affine decomposition. Sections~\ref{sec:pins}
and~\ref{sec:actual-arcs} place plugs on actual rims and prove the bulk
transport theorem. Section~\ref{sec:nondegeneracy} excludes its
degenerate scalar. Section~\ref{sec:observable} constructs the
holomorphic limit and proves Theorem~\ref{thm:cardy}.

Throughout the transport construction, the microscopic limit
$\eps\downarrow0$ is taken first, at a fixed comparison scale $h$ and
for finitely many fixed patterns. Pattern-specific positive
conditioning probabilities may be small. Bounds asserted to be
uniform of order $h$ use only a fixed finite family of short normalized
templates. Spatial refinement depth is also fixed during the
microscopic limit and is subsequently increased before $h$ decreases.
The final Jordan-domain buffers are fixed until the microscopic
crossing limit has been taken. Each argument below specifies which of
these limits is in use.
